%============================================================
% constraints.tex — from the constraints to the algorithm (one-page guide)
% Grow the second table as chapters are written.
%============================================================
\chapter{From Constraints to Algorithm}
\label{ch:constraints}

Greedy problems rarely say ``greedy''. The limits help: $n \le 2\cdot10^5$ with an optimisation question and no obvious DP state is a strong hint that a sort-and-scan greedy is intended. A judge runs roughly $10^8$ simple operations per second.

\section*{Step 1 --- What the limits allow}

\begingroup\small\renewcommand{\arraystretch}{1.3}
\begin{tabularx}{\linewidth}{@{}l l X@{}}
\toprule
\textbf{Limit} & \textbf{Budget} & \textbf{Typical greedy} \\
\midrule
$n \le 8$--$10$ & $\bigO{n!}$ & the brute force you stress-test every greedy against \\
$n \le 5000$ & $\bigO{n^2}$ & lexicographic greedy with an $\bigO{n}$ feasibility check per position (Ch.~\ref{ch:lexicographic}) \\
$n \le 2\cdot10^5$ & $\bigO{n \log n}$ & sort and scan (Ch.~\ref{ch:sorting_greedy}, \ref{ch:intervals}), exchange comparators (Ch.~\ref{ch:exchange}), multisets (Ch.~\ref{ch:containers}) \\
$n \le 10^6$ & $\bigO{n}$ & one pass with a running state or balance (Ch.~\ref{ch:introduction}, \ref{ch:balances}) \\
values up to $10^{18}$ & $\bigO{\log V}$ steps & digit greedy (Ch.~\ref{ch:digits}), reachable sums (Ch.~\ref{ch:math_greedy}) \\
NP-hard problem, large $n$ & polynomial & approximation with a proven ratio (Ch.~\ref{ch:approximation}) \\
\bottomrule
\end{tabularx}
\endgroup

\begin{insightbox}[Test before you prove, prove before you submit]
Write the exhaustive brute force for $n \le 8$ first. Run the greedy against it on thousands of random inputs. A greedy that survives is worth proving; one that fails has just handed you the counter-example.
\end{insightbox}

\clearpage
\section*{Step 2 --- What the question asks}

\begingroup\small\renewcommand{\arraystretch}{1.3}
\begin{tabularx}{\linewidth}{@{}>{\raggedright\arraybackslash}p{6.2cm} X@{}}
\toprule
\textbf{The problem says \ldots} & \textbf{Try first} \\
\midrule
``maximum number of items / people / tasks satisfied'' & take the best available after sorting (Ch.~\ref{ch:introduction}, \ref{ch:sorting_greedy}) \\
intervals, ``meetings'', ``points that hit every interval'' & interval greedy (Ch.~\ref{ch:intervals}) \\
``largest / smallest number after one operation'' & digit greedy (Ch.~\ref{ch:digits}) \\
``lexicographically smallest string / sequence'' & greedy per position, with a feasibility check (Ch.~\ref{ch:strings}, \ref{ch:lexicographic}) \\
workers and jobs, players and trainers & match two sorted lists (Ch.~\ref{ch:matching}) \\
deadlines, ``minimise total waiting / lateness'' & scheduling rules (Ch.~\ref{ch:scheduling}) \\
brackets, fuel, money that must never go negative & running balance (Ch.~\ref{ch:balances}) \\
``every sum from 1 to $n$ must be reachable'' & reachable sums (Ch.~\ref{ch:math_greedy}) \\
``in which order should we do the tasks'' & exchange argument and comparator (Ch.~\ref{ch:exchange}) \\
``take now, maybe undo later'' & greedy with regret (Ch.~\ref{ch:regret}) \\
``assign each request the smallest item that fits'' & sorted containers (Ch.~\ref{ch:containers}) \\
cache, ``evict'', ``future requests unknown'' & online algorithms (Ch.~\ref{ch:online}) \\
weighted intervals, 0/1 knapsack, coin change & greedy fails: the \textit{Dynamic Programming} book (Ch.~\ref{ch:greedy_vs_dp}) \\
heaps of choices, Huffman & the \textit{Stacks, Queues \& Heaps} book \\
minimum spanning tree, shortest paths & the \textit{Graph Algorithms} book \\
\bottomrule
\end{tabularx}
\endgroup

\begin{warningbox}[A plausible greedy is not a proof]
Most wrong answers on greedy problems come from a rule that works on the samples. Every chapter of this book pairs each greedy with a counter-example to a tempting alternative and a written proof.
\end{warningbox}
